% =============================================================================
%  IEEEtran conference manuscript.
%
%  Compiles as-is on Overleaf (select "IEEE Conference" or just upload this
%  file; IEEEtran.cls ships with TeX Live). Locally:
%      pdflatex paper_ieee && pdflatex paper_ieee
%  Two passes are needed for the table cross-references.
%
%  Every numeric result below is taken from the committed artefacts in
%  reports/*.json of github.com/dchaudhari7177/hemophilia and is reproducible
%  from the released scripts.
% =============================================================================
\documentclass[conference]{IEEEtran}
\IEEEoverridecommandlockouts

\usepackage{cite}
\usepackage{amsmath,amssymb}
\usepackage{graphicx}
\usepackage{textcomp}
\usepackage{booktabs}
\usepackage{siunitx}
\usepackage[table]{xcolor}
\usepackage{url}

\sisetup{detect-all, table-number-alignment = right}

\def\BibTeX{{\rm B\kern-.05em{\sc i\kern-.025em b}\kern-.08em T\kern-.1667em\lower.7ex\hbox{E}\kern-.125emX}}

\begin{document}

\title{The Genotype-Only Ceiling for Factor VIII Inhibitor\\Risk Prediction: Patient-Level Modelling and\\Cross-Registry Validation in Haemophilia A}

\author{\IEEEauthorblockN{Dipak Chaudhari, Tejas Nagmote, Sneha A, and Varsha P}
\IEEEauthorblockA{\textit{Department of Computer Science and Engineering} \\
\textit{PES University}\\
Bengaluru, India \\
Project PW\_GRS\_01 \quad Guide: Prof. Gayathri R S}}

\maketitle

\begin{abstract}
Neutralising alloantibodies (``inhibitors'') against infused factor VIII
(FVIII) develop in a substantial minority of people with haemophilia A and are
the most consequential complication of replacement therapy, yet no
risk-stratification tool is in routine pre-treatment clinical use.
Machine-learning studies on public \textit{F8} variant registries report
accuracies of 97\% and above. We show that such figures are not predictive
performance, establish what is genuinely achievable, and determine whether the
limiting factor is the model or the data. Using the EAHAD/HADB \textit{F8}
resource, whose allele-report table contains one row per patient
($n = 10{,}064$) rather than one row per variant, we assembled 4{,}966 records
carrying a recorded inhibitor outcome across 2{,}643 variants and 342
reporting studies. We engineered 95 mechanistic descriptors in six blocks and
screened twelve model families. All partitions were grouped by variant, no
resampling was applied at any stage, and unrecorded outcomes were retained as
unlabelled rather than treated as negatives. A rank-averaged ensemble achieved
AUC-ROC $0.786 \pm 0.004$ under repeated variant-grouped cross-validation. The
genomic-only rung of the feature ablation reached 0.742, independently
reproducing a ceiling previously established on the CDC CHAMP registry;
patient baseline phenotype contributed $+0.023$ AUC. Permuted labels gave
0.504, ungrouped cross-validation inflated the score by 0.032, and a single
outcome-derived column from a widely redistributed derived file moved AUC from
0.777 to 0.966. Cross-registry transfer initially appeared to exceed
within-registry performance; we traced this to 64.5\% variant overlap between
CHAMP and EAHAD, and on genuinely novel variants obtained AUC 0.851 (95\% CI
0.819--0.881) against 0.624 for CHAMP's own cross-validation on identical
rows. Genotype-only prediction saturates near AUC 0.74 in two independently
curated registries; because eleven model families span only 0.027 AUC, this is
a property of the available data rather than of model capacity.
\end{abstract}

\begin{IEEEkeywords}
haemophilia A, factor VIII inhibitors, \textit{F8}, clinical risk prediction,
data leakage, external validation, probability calibration
\end{IEEEkeywords}

% =============================================================================
\section{Introduction}

Haemophilia A is an X-linked bleeding disorder caused by variants in
\textit{F8}, managed by replacement of the deficient coagulation factor VIII.
In a substantial minority of patients the immune system recognises infused
FVIII as foreign and produces neutralising alloantibodies, termed inhibitors.
Once inhibitors develop, replacement therapy loses efficacy, bleeding becomes
markedly harder to control, and the patient requires immune tolerance
induction --- a prolonged and expensive protocol that does not always succeed.

Inhibitor development is concentrated in the first fifty exposure days. This
makes it, in principle, an ideal target for pre-treatment risk stratification:
the window in which the outcome occurs is known and short, and the clinical
response to elevated risk is available and low-cost. Despite this, no
genotype-based stratification tool is in routine clinical use.

The mechanistic basis is well established. Variant classes that abolish
endogenous FVIII protein --- large deletions, nonsense variants and
frameshifts --- carry markedly higher inhibitor incidence than missense
variants, which permit expression of a full-length, if dysfunctional, protein
\cite{gouw}. The immunological interpretation is that a patient who has never
expressed FVIII has established no central tolerance to it, so infused factor
is encountered as a wholly foreign antigen.

Several groups have applied machine learning to public \textit{F8} registries,
most prominently the CDC Hemophilia A Mutation Project (CHAMP) list
\cite{champ}. Reported accuracies reach 97.37\%~\cite{singh}. If reproducible
as \textit{predictive} performance, such figures would be clinically decisive.
They are not reproducible. Prior work by this group established that these
estimates arise from three specific preprocessing choices, and placed the
honest ceiling of genotype-only data near AUC 0.74--0.75. That finding raises
the question this paper addresses: \textbf{is that ceiling a limitation of the
modelling, or of the data?}

CHAMP is a catalogue of \textit{variants}. It records whether an inhibitor has
ever been reported for a variant, but carries no per-patient factor activity,
antigen measurement, or cross-reacting material (CRM) typing. If the ceiling
reflects missing patient-level information rather than insufficient model
capacity, a resource supplying that information should move it.

\subsection{Contributions}
\begin{enumerate}
\item An inhibitor-risk model built at the \textit{patient level} on the
      EAHAD/HADB resource --- one observation per allele report rather than
      per variant --- reaching AUC $0.786 \pm 0.004$ under strict
      variant-grouped validation. We are not aware of prior patient-level
      modelling on this resource but make no priority claim, as no systematic
      review was conducted.
\item Evidence that the genotype-only ceiling of approximately 0.74 replicates
      on an independently curated registry, establishing it as a property of
      the data class rather than an artefact of one dataset or pipeline.
\item Quantification of three leakage mechanisms --- outcome-derived features,
      relabelling of unrecorded outcomes, and ungrouped partitioning --- with
      the AUC attributable to each.
\item Identification and quantification of 64.5\% variant overlap between
      CHAMP and EAHAD, which invalidates naive cross-registry validation in
      this domain, together with a corrected external estimate.
\item Two experiments whose outcomes contradicted our hypotheses, reported
      with the claims revised rather than the framing.
\item A bounded calibration procedure preventing the deployed model from
      emitting risk estimates the underlying cohort cannot support.
\end{enumerate}

% =============================================================================
\section{Related Work and the Reproducibility Problem}

Published analyses of \textit{F8} registries share a common design: rows are
variants, features are registry columns (often label-encoded directly), the
outcome is a reported inhibitor history, and evaluation is a stratified split
or $k$-fold cross-validation. Three properties of this design account for the
reported performance.

\textit{Identifier features.} Columns such as the HGVS cDNA description are
near-unique per row. Label-encoded, they function as primary keys. In our prior
work a random forest given such a column attained training AUC 1.000
\textit{with the outcome labels randomly permuted} --- the diagnostic
signature of an identifier rather than a predictor.

\textit{Relabelling of unrecorded outcomes.} Registry inhibitor fields are
tri-state: reported positive, reported negative, and not reported. Mapping
``not reported'' to negative converts absence of ascertainment into evidence of
absence. In CHAMP this affects 1{,}744 of 4{,}040 rows and reduces apparent
prevalence from 20.1\% to 11.4\%, raising the majority-class baseline against
which accuracy is measured.

\textit{Resampling before partitioning.} Random over-sampling duplicates
minority rows verbatim. Applied before the train/test split, it places
byte-identical copies of training rows into the test set. The reported 97.37\%
accuracy \cite{singh} is obtained under this protocol.

These are not exotic failures. They are the default outcome of applying a
standard tabular machine-learning recipe to a curated biological registry. Our
contribution is to measure each rather than to assert it.

% =============================================================================
\section{Data}

\subsection{The EAHAD/HADB Cohort}
The resource \cite{eahad} is distributed as two supplementary tables at
different levels of analysis: a variant table (6{,}211 rows: mutation class,
protein consequence, FVIII domain, exon, codon and nucleotide change) and an
allele-report table (10{,}064 rows) carrying individual baseline FVIII
activity, clinical severity, antigen, CRM type, reporting centre and source
publication alongside the inhibitor outcome. The latter is the layer absent
from CHAMP, and it determines the design of this study: the unit of analysis
is the patient, with variant annotation joined on the variant identifier.

\subsection{Outcome Definition}
The inhibitor field takes eight surface forms normalising to three states:
4{,}130 negative, 836 positive, and 5{,}098 unrecorded. Unrecorded outcomes
were retained as unlabelled and never mapped to negative. The analytic cohort
is therefore \textbf{4{,}966 records at 16.83\% prevalence}, which falls within
the published range for haemophilia A without adjustment --- itself evidence
that the tri-state field is being read correctly.

\subsection{Cohort Characteristics}

\begin{table}[!t]
\caption{Baseline characteristics of the 4,966 labelled allele reports
(836 inhibitor-positive, 16.8\%), across 2,643 variants and 342 reporting
studies. Median records per variant 1, maximum 104.}
\label{tab:cohort}
\centering
\footnotesize
\begin{tabular}{lrrr}
\toprule
Characteristic & $n$ & \% & Inh.\ rate \\
\midrule
\multicolumn{4}{l}{\textit{Mutation class}} \\
Missense              & 2,677 & 53.9 & 8.3\% \\
Frameshift            & 1,072 & 21.6 & 22.4\% \\
Nonsense              & 594   & 12.0 & 30.1\% \\
Splice                & 250   & 5.0  & 17.2\% \\
Large deletion        & 237   & 4.8  & \textbf{51.9\%} \\
In-frame              & 97    & 2.0  & 27.8\% \\
Silent                & 30    & 0.6  & 3.3\% \\
\addlinespace
\multicolumn{4}{l}{\textit{Clinical severity}} \\
Severe                & 2,718 & 54.7 & 24.1\% \\
Moderate              & 774   & 15.6 & 8.9\% \\
Mild                  & 1,392 & 28.0 & 6.2\% \\
Unknown               & 82    & 1.7  & 29.3\% \\
\addlinespace
\multicolumn{4}{l}{\textit{Baseline FVIII activity (IU/dL)}} \\
$<1$                  & 2,356 & 47.4 & 20.5\% \\
1--2                  & 154   & 3.1  & 8.4\% \\
2--5                  & 610   & 12.3 & 7.4\% \\
5--15                 & 555   & 11.2 & 7.6\% \\
15--40                & 415   & 8.4  & 4.1\% \\
$>40$                 & 30    & 0.6  & 0.0\% \\
Not measured          & 846   & 17.0 & --- \\
\bottomrule
\end{tabular}
\end{table}

Median measured activity was 0.50 IU/dL (IQR 0.50--5.00); 1.00 among
inhibitor-negative and 0.50 among inhibitor-positive records. FVIII antigen was
measured in only 649 records (13.1\%) and CRM type recorded in 1{,}312
(26.4\%). The monotone gradient across mutation classes and across severity
strata (Table~\ref{tab:cohort}) reproduces established immunology and
constitutes a prior check on label quality.

\subsection{The CHAMP Comparison Cohort}
For external validation we used the CDC CHAMP \textit{F8} variant list
\cite{champ}: 4{,}040 variants, of which 2{,}296 carry a recorded outcome at
20.1\% prevalence. CHAMP is variant-level and carries reported severity but no
per-patient measurements.

\subsection{Audit of Redistributed Derived Files}
Two convenience files circulate alongside the raw supplement, prepared by
merging and aggregating the two tables. Both reproduce, on this new dataset,
the failure modes of Section~II. We audit them because they are the form in
which many investigators will encounter this resource.

\textit{Relabelling.} The machine-learning-ready file contains 3{,}706 variant
rows, of which \textbf{1{,}063 carry a negative label} although no inhibitor
outcome was ever recorded for that variant. Apparent prevalence falls from
23.5\% (among variants actually followed up) to 13.3\%.

\textit{Outcome-derived features.} The merged file exposes twelve columns in
our forbidden set, including a per-variant inhibitor positive rate --- the
outcome itself, aggregated over the very patient rows being predicted. Adding
that single column to our design matrix moves AUC from \textbf{0.777 to
0.966}, a change of $+0.189$, with no biological information added.

% =============================================================================
\section{Methods}

\subsection{Feature Construction}
Ninety-five numeric descriptors were derived in six declared blocks: genotype
class (21), FVIII domain (16), protein position (13), residue chemistry (18)
using Grantham distance \cite{grantham} and BLOSUM62 \cite{blosum}, patient
clinical measurements (18), and reporting context (9). Every feature is either
a property of the variant or a measurement recorded at diagnosis; none is
derived from the outcome, and none is row-unique.

Two engineering decisions required judgement. The curated CpG column is empty
in this release, so the hotspot signature was derived from the nucleotide
change, since C\textgreater T and G\textgreater A transitions are the
deamination products of a methylated CpG dinucleotide. The factor-level field
mixes plain numerals, left-censored readings, ranges and annotated entries;
censored values were assigned half the stated bound, the conventional
substitution for a left-censored assay reading, which preserves the ordering of
a censored value below the corresponding observed one. Missingness is encoded
explicitly through indicator features, so median imputation does not silently
create a measurement indistinguishable from an observed one.

\subsection{Validation Protocol}
2{,}643 variants generate 4{,}966 records, with recurrent variants contributing
up to 104 records each. Under random partitioning a model can memorise a
variant during training and be scored on that same variant at test --- the
patient-level form of the identifier leak described in Section~II.
\textbf{Every partition in this study is grouped by variant.}

As a stricter test of transportability we repeat validation grouping by source
publication. Reporting centres differ in inhibitor screening practice and in
which patients they publish, so study-grouped performance estimates behaviour
in a centre not represented in training.

Class imbalance is addressed with class weighting exclusively. No
over-sampling, under-sampling or synthetic minority generation is applied at
any stage.

Because a single 20\% held-out partition contains approximately 159 positive
records and carries roughly $\pm 0.04$ of sampling noise, the primary estimate
is five-fold variant-grouped cross-validation repeated over three seeds on the
entire labelled cohort. A single grouped held-out partition, scored once after
model and thresholds were fixed, serves as an independent check.

\subsection{Models and Ensembling}
Twelve families were screened, spanning linear, kernel, neighbour, forest,
boosting and neural approaches, including XGBoost \cite{xgboost} and LightGBM
\cite{lightgbm}. Leading families were tuned by randomised search over 25 draws
scored on precomputed grouped folds.

The \texttt{scikit-learn} stacking meta-estimator cannot be used in this
setting: it generates internal cross-validation folds at random with no
mechanism to propagate grouping, placing the same variant on both sides of the
meta-learner's training boundary. In our screen this produced an inverted model
(AUC 0.241). We therefore construct the out-of-fold matrix manually from
grouped folds and compare probability averaging, rank averaging, logistic
stacking, and greedy forward selection with replacement \cite{caruana}.

Members are combined by rank average rather than probability average, because
members are calibrated on different scales. Ranks are computed against the
training score distribution stored at fit time rather than within the scoring
batch, so an individual patient's estimate is independent of which other
patients are scored alongside them.

\subsection{Calibration}
Probabilities are calibrated by isotonic regression \cite{isotonic} fitted on
training out-of-fold scores only. Unmodified isotonic regression returned
exactly 1.0 for the highest-scoring inputs, because its terminal step covered a
small, entirely positive training bin. A deployed tool reporting
\textit{certainty} of inhibitor development is not supportable: the
highest-risk observed stratum develops inhibitors at 53.1\%, and the highest
calibration decile at 59.6\%.

We therefore constrain the calibrator. Training scores are partitioned into
deciles, each decile's positive rate is Laplace-smoothed as $(k+1)/(n+2)$, and
predictions are clipped to the interval spanned by those smoothed rates ---
$[0.025, 0.520]$ in this cohort. The transformation is monotone, so ranking and
hence AUC are unchanged; only the reported magnitude is constrained to what the
cohort supports.

\subsection{Statistical Analysis}
Discrimination is reported as AUC-ROC and AUC-PR against the prevalence
baseline. Uncertainty is estimated by stratified bootstrap over 2{,}000
resamples. Paired model comparison uses DeLong's test \cite{delong}.
Calibration is summarised by the Brier score and expected calibration error.
Clinical utility is assessed by decision-curve net benefit \cite{dca}.
Operating thresholds are selected on training folds by Youden's J
\cite{youden}, by accuracy maximisation, and at fixed sensitivity targets.

% =============================================================================
\section{Results}

\subsection{Model Screen}

\begin{table}[!t]
\caption{Twelve model families under five-fold variant-grouped
cross-validation on the training partition ($n = 3{,}992$). The final row
documents a methodological failure and is not a candidate model.}
\label{tab:screen}
\centering
\footnotesize
\begin{tabular}{lrrrr}
\toprule
Model & AUC-ROC & AUC-PR & Bal.\ acc. & Brier \\
\midrule
Extremely rand.\ trees   & 0.7778 & 0.4568 & 0.7183 & 0.1514 \\
Random forest            & 0.7774 & 0.4615 & 0.7176 & 0.1363 \\
LightGBM                 & 0.7745 & 0.4773 & 0.7065 & 0.1420 \\
XGBoost                  & 0.7735 & 0.4678 & 0.7110 & 0.1559 \\
Gradient boosting        & 0.7705 & 0.4769 & 0.7120 & 0.1165 \\
Hist.\ grad.\ boosting   & 0.7682 & 0.4592 & 0.7010 & 0.1479 \\
Multilayer perceptron    & 0.7622 & 0.4361 & 0.7105 & 0.1202 \\
Elastic-net logistic     & 0.7572 & 0.4409 & 0.7002 & 0.1883 \\
L2 logistic              & 0.7552 & 0.4372 & 0.6995 & 0.1887 \\
SVM (RBF)                & 0.7551 & 0.3997 & 0.7018 & 0.1214 \\
$k$-nearest neighbours   & 0.7507 & 0.4078 & 0.6907 & 0.1244 \\
\addlinespace
Stacking (ungrouped CV)  & 0.2412 & 0.1108 & 0.5019 & 0.3492 \\
\bottomrule
\end{tabular}
\end{table}

The eleven valid families in Table~\ref{tab:screen} span 0.027 AUC. Model
capacity is not the binding constraint on this problem --- a finding that
recurs throughout. Tuning improved the leading families to 0.781--0.784. Greedy
forward selection retained three members (XGBoost, extremely randomised trees,
LightGBM) at out-of-fold AUC 0.7859; rank and mean averaging over all five
tuned members gave 0.7860, and the three-member ensemble was retained for
inference economy.

\subsection{Feature Ablation}

\begin{table}[!t]
\caption{Contribution of each feature block under repeated variant-grouped
cross-validation over three seeds.}
\label{tab:ablation}
\centering
\footnotesize
\begin{tabular}{lrrr}
\toprule
Rung & $k$ & AUC-ROC & $\Delta$ \\
\midrule
Genotype class only        & 21 & $0.7024 \pm .0013$ & --- \\
$+$ FVIII domain           & 37 & $0.7261 \pm .0033$ & $+0.0237$ \\
$+$ protein position       & 50 & $0.7372 \pm .0012$ & $+0.0111$ \\
$+$ chemistry (all genomic)& 68 & $\mathbf{0.7417 \pm .0029}$ & $+0.0045$ \\
$+$ patient phenotype      & 86 & $\mathbf{0.7645 \pm .0011}$ & $+0.0228$ \\
$+$ reporting region       & 95 & $0.7779 \pm .0025$ & $+0.0134$ \\
\bottomrule
\end{tabular}
\end{table}

Two findings follow from Table~\ref{tab:ablation}. First, the genomic rung
reaches 0.7417. This reproduces the ceiling previously established on CHAMP ---
a registry curated by a different consortium, from different patients, through
an independently written feature pipeline. Two registries agreeing on where
genotype-only information is exhausted is substantially stronger evidence than
either alone, and reframes the ceiling from a limitation of one analysis into a
property of this class of data. Second, patient baseline phenotype contributes
$+0.023$ AUC, confirming the hypothesis that motivated this study.

\subsection{Leakage Controls}

\begin{table}[!t]
\caption{Controls bounding the primary result.}
\label{tab:controls}
\centering
\footnotesize
\begin{tabular}{lr}
\toprule
Control & Result \\
\midrule
Permuted labels (prior work: 1.000)       & 0.5042 \\
Ungrouped random CV                       & 0.8098 \\
Variant-grouped CV                        & \textbf{0.7779} \\
Study-grouped CV                          & $0.7797 \pm .0004$ \\
Test rows identical to a training row     & 1.03\% \\
Strongest feature--outcome $|r|$          & 0.247 \\
\bottomrule
\end{tabular}
\end{table}

The 0.032 AUC forgone by grouping is of the same magnitude as the entire
clinical-layer gain of 0.023, illustrating how readily a partitioning choice
can be mistaken for a modelling advance. Random over-sampling before
partitioning leaves roughly half the test set byte-identical to training rows;
the observed 1.03\% reflects ordinary coincidence between patients who share a
mutation class, a severity band and the same missing measurements.

\subsection{Primary Result and Operating Characteristics}
Under repeated variant-grouped cross-validation over the full labelled cohort
($n = 4{,}966$), the ensemble achieved \textbf{AUC-ROC $0.7861 \pm 0.0040$}
(per-seed 0.7870, 0.7808, 0.7905).

On the held-out partition ($n = 974$, prevalence 16.3\%): AUC-ROC 0.7451
(95\% CI 0.6927--0.7890); AUC-PR 0.4552 (0.3883--0.5295) against a no-skill
baseline of 0.1632; Brier 0.114; expected calibration error 0.033. The
cross-validated estimate lies within the held-out confidence interval.

\begin{table*}[!t]
\caption{Operating characteristics on the held-out partition ($n = 974$, 159
positive). Majority-class baseline accuracy is 83.68\%. Thresholds were
selected on training folds. The balanced and rule-out points are those at which
the tool is specified; the accuracy-maximising point is reported for comparison
only.}
\label{tab:operating}
\centering
\footnotesize
\begin{tabular}{lrrrrrrrrrrr}
\toprule
Operating point & Threshold & Accuracy & Sensitivity & Specificity & PPV & NPV & MCC & TP & FP & FN & TN \\
\midrule
Accuracy-maximising          & 0.474 & 86.65\% & 25.8\% & 98.5\% & 77.4\% & 87.2\% & 0.396 & 41  & 12  & 118 & 803 \\
Balanced (Youden's J)        & 0.181 & 76.69\% & 61.6\% & 79.6\% & 37.1\% & 91.4\% & 0.343 & 98  & 166 & 61  & 649 \\
Rule-out (80\% sens.\ target)& 0.131 & 65.61\% & 74.2\% & 63.9\% & 28.6\% & 92.7\% & 0.285 & 118 & 294 & 41  & 521 \\
Rule-out (90\% sens.\ target)& 0.083 & 49.38\% & 79.9\% & 43.4\% & 21.6\% & 91.7\% & 0.176 & 127 & 461 & 32  & 354 \\
\bottomrule
\end{tabular}
\end{table*}

Accuracy is reported only alongside its baseline. The accuracy-maximising
threshold in Table~\ref{tab:operating} exceeds the no-skill baseline by 2.97
percentage points while identifying 25.8\% of inhibitor-positive patients; it
is the least clinically useful row in the table, and the one that an evaluation
rubric specifying ``high accuracy'' would select. The tool is specified at the
balanced and rule-out points, where negative predictive value exceeds 0.91.

Decision-curve net benefit was 0.088 at a 10\% threshold probability (treat-all
0.070) and 0.058 at 20\% (treat-all $-0.046$). The model exceeded both
treat-all and treat-none strategies for threshold probabilities between 0.09
and 0.52, covering 73\% of the swept range. Below a 9\% threshold, monitoring
every patient is preferable and the model offers no advantage.

Against a genomic-only model on the same held-out patients (AUC 0.7305), the
full model's advantage of 0.0147 did not reach significance (DeLong $z = 1.06$,
$p = 0.29$), which is expected at this sample size. The ablation in
Table~\ref{tab:ablation}, estimated over the full cohort with repeated
resampling, is the better-powered comparison.

\subsection{Cross-Registry Validation and a Registry-Overlap Artefact}
Projecting both registries into a 57-feature intersection, a model trained on
HADB and applied to CHAMP without refitting scored AUC 0.879 --- exceeding
CHAMP's own within-registry cross-validation of 0.725.

A model applied to an unseen cohort should not outperform that cohort's own
cross-validation. We treated this as a diagnostic rather than a result. CHAMP
and EAHAD are both curated from the published literature, and they curate
substantially the same publications. Matching on mutation class, mature
residue, reference residue and substituted residue, \textbf{1{,}259 of 1{,}951
(64.5\%)} of CHAMP's substitution-resolvable labelled variants are already
present in HADB.

One methodological subtlety governs this comparison. The matching key requires
a reference residue and position, which frameshift and large structural
variants do not provide. Left uncorrected, every frameshift is classified as
novel by default, and the resulting comparison contrasts a 77\%-frameshift
subset with a 60\%-missense subset, measuring mutation-class composition rather
than novelty. Both strata are therefore restricted to substitution-like
variants.

\begin{table}[!t]
\caption{Transfer from HADB to CHAMP, stratified by prior presence of the
variant in the training registry.}
\label{tab:transfer}
\centering
\footnotesize
\begin{tabular}{lrrr}
\toprule
Stratum & $n$ & Prev. & AUC-ROC \\
\midrule
Present in HADB            & 1,259 & 0.137 & 0.9361 \\
Novel to HADB              & 692   & 0.240 & \textbf{0.8513} \\
Attributable contamination & ---   & ---   & 0.0848 \\
\bottomrule
\end{tabular}
\end{table}

The corrected external estimate is AUC 0.8513 (95\% CI 0.8193--0.8807). On the
identical 692 novel records, CHAMP's own cross-validated model achieved AUC
0.6237. Training on 4{,}966 patient-level European records therefore predicted
novel North American variants \textbf{0.2276 AUC better} than CHAMP predicted
them from its own data --- independent confirmation, by a different route, of
the ablation's conclusion that repeated patient-level observation constitutes
richer supervision than a variant catalogue.

\subsection{Reproduction of the Relabelling Effect}

\begin{table}[!t]
\caption{Identical fitted model and variant set; only the label convention
differs.}
\label{tab:relabel}
\centering
\footnotesize
\begin{tabular}{lrrrrr}
\toprule
Label convention & $n$ & Prev. & Base & Acc. & AUC \\
\midrule
Recorded outcomes only        & 2,296 & 0.201 & 79.9\% & 80.1\% & 0.879 \\
Unrecorded $\rightarrow$ neg. & 4,040 & 0.114 & \textbf{88.6\%} & 74.7\% & 0.855 \\
\bottomrule
\end{tabular}
\end{table}

Discrimination does not improve under relabelling; AUC in fact declines. What
changes is the accuracy obtainable \textit{without skill}, which rises from
79.9\% to 88.6\%. A headline accuracy near 97\% is measured on a set whose
no-skill baseline is already 88.6\%, and is approached by tuning the decision
threshold toward the enlarged majority class.

\subsection{Results Contrary to Hypothesis}
We predicted that collapsing records to one row per variant would lose
information. It scored \textit{higher}: 0.8038 against 0.7766. The two
estimates are not comparable, differing in unit of analysis, label definition
and sample size, and predicting a variant's \textit{modal} outcome is an easier
problem than predicting an individual's, since feature averaging across a
variant's records cancels measurement noise and well-characterised recurrent
variants dominate the row set.

The case for patient-level modelling therefore rests on a different
measurement. Among the 537 variants with two or more recorded outcomes,
\textbf{124 (23.1\%) are discordant} --- patients carrying an identical variant
who differ in inhibitor outcome --- spanning 1{,}297 records. Aggregation
assigns every such patient the same prediction and is necessarily wrong for the
minority, and it cannot use per-patient factor activity at all. The second
contrary result, naive cross-registry transfer overstating external validity,
is detailed in Section~V-E.

\subsection{Feature Attribution}
Permutation importance on the held-out partition ranked severity ordinal
(0.0084 AUC), East Asian reporting region (0.0073), nucleotide position
(0.0046), log FVIII activity (0.0044) and multi-domain involvement (0.0043)
highest. Aggregated by block: clinical 0.0287, position 0.0211, context 0.0148,
chemistry 0.0116, domain 0.0088, genotype 0.0042.

% =============================================================================
\section{Discussion}

\subsection{What Limits Performance}
The evidence converges on a single conclusion: the constraint is
informational, not architectural. Eleven model families span 0.027 AUC; tuning
contributed 0.005 and ensembling 0.003. By comparison, the partitioning leak we
declined to exploit was worth 0.032, and a single outcome-derived column was
worth 0.189.

The genomic ceiling of approximately 0.742, now observed in two independently
curated registries, is best interpreted as the information content of
\textit{F8} genotype with respect to inhibitor development. Established
determinants that no public registry records --- treatment intensity, exposure
days at first bleed, product class (plasma-derived versus recombinant
\cite{sippet}), surgery or infection at first exposure, and HLA class II typing
--- are absent by construction. Their absence, not model capacity, bounds these
estimates.

\subsection{Direction of Effect for Clinical Features}
Baseline factor activity could in principle be a \textit{consequence} of the
outcome rather than a predictor of it, since a circulating inhibitor suppresses
measured FVIII activity. The registry documents these as diagnostic baseline
values but does not timestamp them relative to inhibitor detection, so this
concern cannot be fully resolved with these data.

We therefore report the genomic-only rung alongside every headline. A reader
who rejects the clinical block entirely retains a defensible 0.742 model
validated across two registries. The clinical contribution is presented as an
improvement conditional on that assumption, not as a foundation.

\subsection{The Reporting-Region Block}
The context block contributes $+0.0134$ AUC, and East Asian reporting region
ranks second in permutation importance. This warrants explicit caution. The
registry records the reporting laboratory's country, not patient ancestry.
Observed regional inhibitor rates vary widely (East Asia 45.5\%, $n = 308$;
South Asia 10.8\%, $n = 297$; North America 90.9\%, $n = 11$), and the extreme
values rest on very small strata.

Ancestry is a genuine, established risk modifier for inhibitor development, so
part of this signal may be biological. But it is confounded with centre-level
screening intensity and publication selection, and we cannot separate the two
here. A deployment unwilling to accept that confound can use the model without
the context block at AUC 0.7645, and we recommend this for any use outside the
regions represented in training.

\subsection{An Anomaly in CRM Typing}
CRM type did not behave as the underlying immunology predicts. CRM-negative
(type I) patients, those with no circulating antigen, showed a \textit{lower}
inhibitor rate (3.8\%, $n = 156$) than CRM-unclassified patients (22.6\%,
$n = 953$). We attribute this to ascertainment: CRM typing is performed
selectively, largely on mild and moderate patients with discrepant one-stage
and chromogenic assays, who are at low baseline inhibitor risk. This is a
selection effect in the registry rather than a contradiction of the
immunology, but it means CRM type should not be interpreted causally in this
cohort.

\subsection{Implications for Cross-Registry Validation}
The 64.5\% overlap between CHAMP and EAHAD has consequences beyond this study.
Cross-registry transfer is widely regarded as strong evidence of external
validity, and in this domain a naive application of it inflated our own
estimate by 0.085 AUC. Investigators validating \textit{F8} models --- and
plausibly models built on other curated disease-variant resources compiled from
shared literature --- should establish variant-level disjointness before
reporting external performance, and should report the matching procedure, since
key resolvability interacts with mutation class in a way that can silently
confound the comparison.

\subsection{Clinical Positioning}
The model is a research tool and is not a validated medical device. It does not
replace laboratory inhibitor testing or clinical judgement. Its intended use is
prioritisation of monitoring intensity during the first fifty exposure days. At
the balanced operating point it identifies 61.6\% of patients who go on to
develop inhibitors, at a cost of 166 false positives per 974 patients, with
negative predictive value 0.914. The clinical consequence of a false positive
is intensified monitoring, not a change in diagnosis or a withheld therapy ---
the asymmetry that makes a rule-out threshold defensible at this level of
discrimination.

% =============================================================================
\section{Limitations}
\begin{enumerate}
\item Both cohorts are compiled from published reports rather than consecutive
      series. Patients with inhibitors are plausibly more likely to be
      published, and screening intensity varies by centre.
\item The temporality of baseline factor measurements relative to inhibitor
      detection is unresolved (Section~VI-B).
\item Causal covariates --- treatment intensity, exposure days, product class,
      HLA typing --- are absent (Section~VI-A).
\item Reporting region is a confounded proxy for ancestry (Section~VI-C).
\item The held-out partition is small ($n = 974$, 159 positive), giving wide
      intervals; the repeated cross-validated estimate is the primary result
      for this reason.
\item Scope is a single gene and a single disease; transfer to haemophilia B
      (\textit{F9}) was not attempted.
\item All estimates are retrospective and registry-based. No prospective
      validation was performed.
\end{enumerate}

% =============================================================================
\section{Conclusion}
Genotype-only prediction of FVIII inhibitor risk saturates near AUC 0.74, a
value we observe independently in two separately curated registries and which
is therefore best understood as a property of the available data rather than of
model capacity. Patient-level records carrying individual baseline phenotype
raise achievable discrimination to $0.786 \pm 0.004$ under strict
variant-grouped validation, with calibrated probabilities and retained
performance in unseen reporting centres.

Published accuracies above 95\% do not represent predictive performance. Each
contributing mechanism --- outcome-derived features, relabelling of unrecorded
outcomes, and resampling before partitioning --- is reproducible on demand and
quantified here.

Progress beyond approximately 0.79 will not come from better classifiers. It
requires prospective cohorts recording treatment intensity, exposure history,
product class and HLA type. Until such data are available, the honest
contribution of genotype-based modelling is a calibrated rule-out instrument,
not a decisive diagnostic.

\section*{Data and Code Availability}
All code, trained artefacts, result files and the executed analysis notebook
are available at \url{https://github.com/dchaudhari7177/hemophilia}. Every
result is regenerated by the released scripts and serialised to JSON. The
analysis notebook executes eight integrity assertions, including a direct test
for the over-sampling signature, and fails if any does not hold. The test suite
comprises 128 tests.

% =============================================================================
\begin{thebibliography}{00}

% ----------------------------------------------------------------------------
% VERIFY BEFORE SUBMISSION. The two entries below are described rather than
% reconstructed from memory, so that no incorrect citation is propagated.
% ----------------------------------------------------------------------------
\bibitem{singh} Singh and Singh, ``Random-forest classification of \textit{F8}
variants for inhibitor prediction,'' 2025. Reports 97.37\% accuracy.
\textsc{[Citation to be completed from the primary source.]}

\bibitem{eahad} EAHAD/HADB coagulation factor VIII variant resource,
\textit{Blood Advances}, supplement VTH-2024-000215, 2024.
\textsc{[Author list, volume, pages and DOI to be completed.]}
% ----------------------------------------------------------------------------

\bibitem{champ} A. B. Payne, C. H. Miller, F. M. Kelly, J. M. Soucie, and
W. Craig Hooper, ``The CDC Hemophilia A Mutation Project (CHAMP) mutation
list: a new online resource,'' \textit{Human Mutation}, vol. 34, no. 2,
pp. E2382--E2391, 2013.

\bibitem{grantham} R. Grantham, ``Amino acid difference formula to help explain
protein evolution,'' \textit{Science}, vol. 185, no. 4154, pp. 862--864, 1974.

\bibitem{blosum} S. Henikoff and J. G. Henikoff, ``Amino acid substitution
matrices from protein blocks,'' \textit{Proc. Natl. Acad. Sci. USA}, vol. 89,
no. 22, pp. 10915--10919, 1992.

\bibitem{xgboost} T. Chen and C. Guestrin, ``XGBoost: a scalable tree boosting
system,'' in \textit{Proc. 22nd ACM SIGKDD Int. Conf. Knowledge Discovery and
Data Mining}, 2016, pp. 785--794.

\bibitem{lightgbm} G. Ke et al., ``LightGBM: a highly efficient gradient
boosting decision tree,'' in \textit{Advances in Neural Information Processing
Systems}, 2017.

\bibitem{caruana} R. Caruana, A. Niculescu-Mizil, G. Crew, and A. Ksikes,
``Ensemble selection from libraries of models,'' in \textit{Proc. 21st Int.
Conf. Machine Learning}, 2004.

\bibitem{isotonic} B. Zadrozny and C. Elkan, ``Transforming classifier scores
into accurate multiclass probability estimates,'' in \textit{Proc. 8th ACM
SIGKDD Int. Conf. Knowledge Discovery and Data Mining}, 2002, pp. 694--699.

\bibitem{delong} E. R. DeLong, D. M. DeLong, and D. L. Clarke-Pearson,
``Comparing the areas under two or more correlated receiver operating
characteristic curves: a nonparametric approach,'' \textit{Biometrics},
vol. 44, no. 3, pp. 837--845, 1988.

\bibitem{dca} A. J. Vickers and E. B. Elkin, ``Decision curve analysis: a novel
method for evaluating prediction models,'' \textit{Medical Decision Making},
vol. 26, no. 6, pp. 565--574, 2006.

\bibitem{youden} W. J. Youden, ``Index for rating diagnostic tests,''
\textit{Cancer}, vol. 3, no. 1, pp. 32--35, 1950.

\bibitem{sippet} F. Peyvandi et al., ``A randomized trial of factor VIII and
neutralizing antibodies in hemophilia A,'' \textit{New England Journal of
Medicine}, vol. 374, no. 21, pp. 2054--2064, 2016.

\bibitem{gouw} S. C. Gouw, H. M. van den Berg, J. Oldenburg et al.,
``\textit{F8} gene mutation type and inhibitor development in patients with
severe hemophilia A: systematic review and meta-analysis,'' \textit{Blood},
vol. 119, no. 12, pp. 2922--2934, 2012.
\textsc{[Page range to be verified.]}

\end{thebibliography}

\end{document}
